% Part: lambda-calculus
% Chapter: lambda-definability
% Section: lists

\documentclass[../../../include/open-logic-section]{subfiles}

\begin{document}

\olfileid[pt]{lam}{ldf}{lst}
\olsection{Listas}

A lista $[M_1, M_2, \ldots, M_n]$ pode ser definida como:
\[
  [M_1, M_2, \ldots, M_n] \ident \lambd[sz][s M_1 (s M_2 (\ldots (s M_n z)))]
\]
Em termos informais, o termo que representa uma lista é a função
``acumuladora'' da lista: uma função que aceita dois termos; o primeiro
é uma função que recebe o valor acumulado junto com um novo elemento e
retorna o novo valor acumulado; o segundo é o valor acumulado inicial;
ela acumula os elementos um a um e, por fim, retorna o resultado.

\begin{digress}
  Se você conhece programação funcional, verá que a definição é semelhante
  à \emph{redução à direita}: uma lista é definida como a função de
  redução à direita dos elementos que contém.
\end{digress}

Podemos definir facilmente algumas funções úteis com base nessa codificação; 
\begin{align*}
  \fn{Sum} & \ident
  \lambd[l][l \, (\lambd[xa][\fn{Add} \, x \, a])\,  \num{0}]\\
  \fn{Len} & \ident
  \lambd[l][l \, (\lambd[xa][\fn{Add} \, a \, \num{1}]) \num{0}]
\end{align*}

$\fn{Sum}$ calcula a soma de uma lista de numerais de Church. Seu funcionamento
consiste em realizar uma acumulação sobre a lista, na qual o valor inicial é
$\num{0}$ e, para cada elemento, toma-se $\fn{Add} \, x\, a$ como o novo
valor (em que $x$ é o elemento atual e $a$ é o valor acumulado). O resultado
é a soma de todos os elementos.

\begin{prob}
  O que $\fn{Len}$ faz? Explique.
\end{prob}

\begin{prob}
  Defina uma função que inverta uma lista.
\end{prob}
\end{document}

